\chapter[Printing, Readers and Arguments]{How Printing Made More Readers and More Arguments}
\section{An Unsigned Leaflet}
Pick up no treasure or antique, just the printed leaflet you might receive outside a station or market or find on a school noticeboard.

Thin, printed on one side with evenly sharp letters, it rarely earns a second glance. Notice its sameness: thousands of copies carry identical characters, not anyone's handwriting. Behind them stands a machine, not a visible writer. The sheet costs so little it can be given away; your willingness to read is precious.

Take it six centuries back to a medieval European and it seems magical. Books were \textbf{copied} individually. A skilled scribe might spend years on a Bible, using skins from hundreds of sheep, costing years of ordinary food. Books were family assets and sacred objects, not commodities or casual reading. A village priest might alone have access to one, chained beside a table near the altar against theft.

Your nearly worthless sheet can become a thousand identical copies overnight, carried by a thousand hands to a thousand streets.

Our protagonist is neither author nor book but the ability to \textbf{copy} text: printing. What happens when words become cheap? More readers and arguments, wider knowledge and rumours, liberation and amplified hatred. The question remains open as your life inhabits a still fiercer copying revolution.

\section{Market Day in 1520}
Enter a scene.

An autumn market day in a small German town, 1520. Rain fills paving gaps; livestock, baking bread, and fresh, pungent ink scent the air.

Cloth, pot, and herring sellers fill the square, but a travelling bookseller draws the biggest crowd. No great volumes, only pamphlets of a few to dozens of pages, fastened into booklets affordable for a day or two of an artisan's wages.

More than half concern a monk, Martin Luther, quarrelling with Rome.

Three years earlier, he posted ninety-five propositions at Wittenberg, especially opposing indulgences, certificates supposedly reducing posthumous punishment for payment. Written in Latin for academic debate, they had perhaps only dozens of local readers. Under older conditions, scholars would have argued and the matter faded.

Instead, someone printed them, then translated and printed German versions. Within weeks, material for a few scholars reached towns across German lands. Now dozens of Luther's pamphlets sit on this stall, sold like the year's favourite songs.

A literate carpenter reads aloud; his illiterate wife listens intently at the front. A priest scowls outside the circle, holding an authorised rebuttal. Church and critics share machines, paper, and stall. Nearby buyers whisper together. Charles V's order forbidding printing, sale, and reading of Luther hangs on the wall, paste barely dry while business continues.

Remember ink, carpenter, priest, and ignored prohibition. The rest of the chapter asks what is happening in this square.

\section{Books Are Cheaper. What Next?}
State the conflict before answers.

Everyone judges printing differently, plausibly.

The bookseller celebrates commerce and access: once house-priced books now meal-priced pamphlets, available to people previously excluded. Surely knowledge leaving courts and monasteries is wonderful?

Hear the apparently outdated priest before laughing. Trained writers, copyists, and custodians once filtered books. Now anyone prints anything: Luther today, lies and incitement tomorrow. Who checks a leaflet reaching ten cities overnight? High barriers also filtered; removing them admits mud with clean water.

The carpenter's wife asks more simply: what do countless pamphlets offer if I still need someone to read? How can I verify his reading? Does the reader become another priest?

Together they ask:

\textbf{When copying costs collapse, do we gain more truth or more voices? Are these identical?}

If identical, cheaper expression should automatically make humanity wiser and freer. Otherwise we must identify the mixture of voices, replacement gatekeepers, and readers needed when everyone can author.

The 1520 question still burns. If you ruled then, would you ban the flying pamphlets?

\section{Printers, Censors, and Excluded Readers}
Strengthen each case before looking beyond Europe, since printing began earlier elsewhere.

\textbf{First, state the censor's case fully.}

Calling prohibitions reactionaries' fear of truth is not entirely wrong but too easy. Their strongest reasons remain in use today.

First, misinformation kills. Pamphlets spread poisoned-well accusations against Jews, followed by violence. Ten thousand printed lies become ten thousand authoritative-looking statements, not private whispers. Second, public theological disagreement hardens into camps and war, unlike negotiable episcopal discussion. Subsequent religious wars killing millions vindicated that fear. Third, gatekeepers sincerely believed dangerous books condemned souls. Within that worldview, censorship resembled a parent keeping hazardous instructions from children.

This does not excuse censorship. \textbf{Its impulse arises not from stupidity but genuine fear of losing control}, persisting under new technologies.

\textbf{Now those wanting to print.}

Printers were merchants and early media workers, covertly welcoming prohibitions that advertised forbidden books while risking confiscation. Authors discovered unprecedented reach to thousands of strangers. Luther later praised printing as God's gift spreading the gospel. Bypassing intermediaries enabled revolutionary and dangerously unconstrained influence.

\textbf{Finally, those unable to read.}

Early printing generations left most Europeans illiterate, with even fewer literate elsewhere. More readers first amplified \textbf{literate urban people who could afford books}: artisans, townspeople, small traders. Peasants, women, and poor people largely depended on listening. The new public square still charged literacy and money for admission, a ticket we will revisit.

Now leave Europe or misrepresent the story.

\textbf{East Asia printed earlier and differently.} Tang China used carved woodblocks from roughly the seventh to ninth centuries, reversing whole pages on wood for hundreds or thousands of impressions. The earliest clearly dated surviving woodblock print is the 868 Diamond Sutra from Dunhuang, its illustrated opening already technically mature, nearly six centuries before Gutenberg. Northern Song Bi Sheng developed clay movable type; Goryeo Korea advanced metal type, with the surviving Jikji printed in 1377, over seventy years before Gutenberg's Bible.

Why then does printing revolution usually mean Europe? Distinguish facts first. Thousands of Chinese characters complicated movable type, whereas Europe's few dozen letters lowered barriers. Reusable woodblocks remained practical in China. Late-Ming Jiangnan publishers filled markets with novels, drama, and daily encyclopaedias; Korea and Japan also flourished in publishing. Both worlds changed through different combinations: European alphabets, religious fractures, and commerce; Chinese exams, scholarly culture, and broad popular reading.

\textbf{A reverse example} comes from Ottoman lands, where Arabic-script printing faced long resistance and restriction, with Istanbul's first secular Arabic-script press appearing only in 1727. Copyist guilds, religious caution, and manuscripts' artistic and social status mattered. \textbf{Invention does not automatically conquer society}: people embrace, modify, or restrain technologies. Technology is seed; society soil.

\section[Thinking Bridge: The Cost of Copying]{Thinking Bridge: Calculate the Copying Accounts}
Instead of merely marvelling at printing, use book historians' portable accounts, separating several calculable effects.

\textbf{First account: copying costs.}

Hand-copying a thousand texts requires a thousand acts, costs broadly proportional to copies. Printing front-loads expense into the \textbf{first setting}, then adds little per copy. Per-copy economics becomes per-edition economics, favouring expected large sales and \textbf{popular taste}. A specialist manuscript sustained by one patron may never reach type, while sensational pamphlets become profitable.

\textbf{Second account: copying precision.}

Scribes introduce errors, omissions, and personal improvements; copies differ. Printing makes thousands identical, enabling stable references to page five, line three. But \textbf{errors are standardised too}. Local manuscript variation becomes mass identical error with printed authority. Precision neutrally amplifies whatever enters.

\textbf{Third account: gatekeeper location.}

Every system has selectors. In manuscript culture, gatekeepers control the \textbf{entrance}: writing, copying, and custody. Printing widens access and pushes control towards the \textbf{middle and exit}, through banned-book lists, licences, burning, and pursuit. Gatekeepers move, not vanish, awkwardly chasing mobile copies past prohibitions. Who controls what others read?

Combine the accounts: when technology supposedly frees information, ask which content becomes profitable, what besides truth gets replicated, and where new gatekeepers hide.

These questions work for print, newspapers, radio, television, and internet. That is transferability.

\section{Shen Kuo's Piece of Clay}
Read one of humanity's earliest surviving movable-type accounts, in a Chinese official's notebook.

Northern Song Shen Kuo studied astronomy, mathematics, geology, medicine, and music. His later Dream Pool Essays recorded many things, including an artisan he had heard about:

\begin{quote}
During Qingli, a commoner named Bi Sheng made movable plates. He carved characters in clay, as thin as a coin's edge, one stamp per character, and fired them hard. (Dream Pool Essays, volume 18, Techniques.)
\end{quote}
Qingli ran from 1041 to 1048. Shen explains setting fired clay characters on an iron plate coated with pine resin, wax, and paper ash. Heating fixed the plate; reheating released reusable type. Wood had been rejected because variable grain swelled unevenly in ink and stuck to resin.

Notice the irony: history's first named movable-type inventor remains scarcely known, without secure birth or death dates, identified publications, or surviving work. \textbf{His whole recorded existence occupies roughly two hundred characters in a scholar's notes}. Shen recorded curiosity, not an official honour. A world-changing skill in ordinary hands needed an attentive person's casual note to enter history.

This corrects Gutenberg invented printing. More precisely, he independently combined metal type, oil-based ink, and a screw press into an efficient system: integration rather than creation from nothing. Bi Sheng preceded him by four centuries, the Diamond Sutra by six, Jikji by seventy years. A lone German genius lighting medieval darkness wrongs East Asia and obscures why similar technologies produce different social fruits.

\section{Was the Press Protestantism's Engine?}
Return to 1520. Separate wide circulation and Reformation as events, causal explanation, participants' divine-gift or devil's-trumpet accounts, and our judgement. Ask \textbf{how printing and Reformation were causally related}.

\textbf{Explanation one: printing drove reform.}

Elizabeth Eisenstein developed this influential case: without print, Luther disappears into another scholarly quarrel. Print rapidly spread his ideas and vernacular Bibles, allowing direct reading and judgement beyond priests and breaking interpretive monopoly. Timing and mechanism fit: reform followed widespread printing within seventy years. Yet Catholic Spain and Italy printed extensively too, while eight hundred years of East Asian printing produced no corresponding Reformation. Why did the engine ignite only some places?

\textbf{Explanation two: amplifier, not engine.}

Corruption, indulgence grievances, princes' resentment, and new urban groups already cracked the church's order. Print spread emerging lava faster and further. Crucially, \textbf{the church printed too}: indulgences themselves were mass prints, and rebuttals filled carts. Both sides fought the first printing war at shared stalls. Social conflict leads; technology catalyses. But catalysts may alter reactions: unamplified complaints might remain local and negotiable, as earlier reforms were absorbed. Sufficient quantity changes quality.

\textbf{Explanation three: print changed the ecology of opinion.}

Rather than which side won, ask how the \textbf{rules} changed. Disagreement became \textbf{public, lasting, and cumulative}. Manuscript heresies were small fires readily stamped out; thousands of orthodox, heretical, moderate, and extreme copies could not all be recovered or burned. Backed-up opinions demanded debate rather than eradication. Public life entered unstoppable argument, not guaranteed truth. Different social permission for print networks explains regional divergence and later developments. Yet this broad ecology cannot decide whether people became wiser or more dogmatic.

Each holds truth. Omnipotent technology is arrogance; irrelevant technology is obtuseness. Existing fractures plus irreversible amplification better approximate history.

Irreversibility soon reveals a frightening side.

\section[Further Challenge: Reading Nations]{Further Challenge: Nations Read into Being}
Our weightiest theory asks where \textbf{French and German peoples} came from.

A medieval Provençal and Norman farmer might never meet, understand each other's speech, or know each other's lives. Why compatriots? Village and God might be their only belongings. How did nationhood enter strangers' minds strongly enough for mutual sacrifice?

Benedict Anderson's 1983 Imagined Communities answers with the \textbf{imagined community}. Imagined means not false but membership among mostly unknown companions, made possible through print.

Follow the links. Profit requires markets beyond limited Latin readers, hence vernacular German, French, and English. Too many dialects cannot all be printed, so publishers favour urban varieties, gradually smoothing them into written standards. Then northern and southern Germans read the same news and pamphlets, reacting together without meeting. \textbf{Unseen companions reading alike produce a printed we}. Newspapers intensify the daily shared ritual, nationhood's morning prayer.

Media become community builders, not merely channels: reading helps form groups rather than groups merely finding reading matter. Multilingual Ottoman and Qing empires face language-bounded publics increasingly claiming self-rule, with print behind modern imperial breakup and nation-states.

But explanation is not endorsement. It accounts for solidarity and hostility, since we implies they. Print is not the only loom: schools, military service, railways, maps, and censuses also matter, as Anderson acknowledged. Nor should a theory developed around Europe and Latin America transfer carelessly to East Asia: shared Chinese writing sustained a vast imagined community, joining dialects where Europe divided Latin's world. Same machine, opposite fruits.

\section{Everyone Is a Printer}
Before today, counter the naive belief that \textbf{more printed books automatically bring reason and freedom}.

The 1486 Malleus Maleficarum, less than forty years into European printing, taught witch-identification, interrogation, and punishment under two inquisitors' names. Repeated editions made this persecution manual a reference for nearly two centuries of witch-hunting, killing thousands of women and some men. Its exact toll is uncountable, but print armed persecution as efficiently as truth. Anti-Jewish pamphlets repeatedly renewed accusations such as child murder. Luther, witch manuals, and hate all entered homes through the same unselective machine.

Turn the accounts towards today.

Internet-print comparisons often stop at cheaper distribution. Our tool goes further.

Costs: printing reduced them by orders of magnitude; internet copying reaches zero, even without first-edition investment. Gutenberg needed presses and type; your post needs no upfront outlay. \textbf{Everyone is a printer}. Expected sales no longer decide worthiness to circulate. Attention, not copying cost, constrains content, rewarding emotion over accuracy. Popular taste becomes intensified algorithmic extremes.

Precision: one printed error reached ten thousand; a rumour now reaches ten million within an hour, endorsed by friends. Screenshots, reposts, and archives prevent deletion. Backups that protected heresy now protect misinformation. The mechanism serves both good and harm.

Gatekeepers seem gone because everyone can publish, but hide in recommendations, trending lists, and moderation. Code optimised to retain you selects the feed rather than you or an editor. Old banned-book lists publicly identified prohibitions; opaque algorithms may leave you unable to explain repeated viewpoints. \textbf{Visible gatekeepers can be opposed; invisible ones conceal even the target of opposition}. Charles V could scarcely imagine it.

History does not simply repeat. Printing also fostered journals and peer review, libraries, journalistic professionalism, and compulsory schooling: \textbf{new filters built after society experienced loss of control}. Our internet-era filters, fact-checking, media education, platform accountability, remain under construction. The priest's mud did arrive; history's response was new filters, not restoring an irretrievable old one.

\section{Your Question: Who Keeps the Gate?}
Return to the square.

Imperial prohibitions hang while pamphlets sell. Our earlier ban-or-not question now becomes:

\textbf{When copying and transmission approach zero cost, who should be the gatekeeper?}

Answer with reasons after weighing the evidence.

For gatekeepers: rumours kill, manuals and hate are real, the least able to judge drown first, and selection never disappears, only becomes public or algorithmically hidden.

Against them: almost every historical gatekeeper abused power, suppressing later truths alongside lies. Surrendered filtering authority is hard to recover, and interested parties define danger, as questioning indulgences once counted dangerous.

Control may also fail despite intention: pamphlets pass decrees and backups survive bonfires. Yet an unmanaged square mixes water and faster-flowing mud.

With infallible gatekeepers nonexistent, must we choose dangerous freedom or more dangerous order? Could readers instead become their own gatekeepers? Then literacy and education quietly charge admission again.

There is no standard answer. This chapter too is a printed, author-, editor-, and publisher-filtered object, trusted partly through invisible filters. When the next screen message angers or excites you, ask which filters it passed through or bypassed to arrive.
